\section{统一角度--频率--偏振响应表述}

\subsection{以响应场作为首要设计对象}

设 $\G$ 表示一个满足工艺约束的自由形态二值超表面图案。与其用单个性能指标来概括器件行为，我们为每个结构定义如下响应张量：
\[
\mathbf{R}(\G)=\left\{\tp(\lam_i,\ang_j),\ts(\lam_i,\ang_j)\right\}_{i=1,\ldots,N_{\lambda},\,j=1,\ldots,N_{\theta}},
\]
其中两个通道可在需要时推广到更一般的偏振基。该对象记录了在有限角度--频率采样窗口上的算子相关传输特性，因此比单点效率更符合傅里叶算子超表面的物理本质。

这种以响应场为中心的定义对于傅里叶算子器件尤为自然。器件功能本质上由空间频率域中的传递函数形状决定，而在具体照明和成像几何下，这一空间频率变量与入射角直接耦合。一旦采用这一表述，传统上被视为不同器件类别的“微分器”“滤波器”或“偏振选择性算子”，便都可看作同一响应空间中的不同目标几何。

\subsection{统一目标家族}

在所提出的响应空间视角下，多种看似不同的傅里叶算子器件可以用同一种数学语言表达。二阶微分器可写成一个在角向上具有特定偶对称形状的目标传递轮廓，并可扩展到多波长条件。偏振复用器件则对应于为不同偏振通道指定不同的目标响应。光谱选择性算子则体现为同一目标轮廓族随波长系统变化。尽管这些物理行为不同，但逆向变量始终是 $\G$，目标始终是 $\Rtar(\lam,\ang,p)$。

这一统一并非仅仅便于组织章节。更重要的是，它说明生成模型的真正任务不是在有限几个预定义应用之间做分类，而是学习“给定目标响应场时，对应结构分布”这一条件概率。不同任务外观上的差异，被吸收到同一响应流形上的不同目标切片之中。

\subsection{物理一致的逆向目标}

我们希望找到一个或多个结构，使其验证响应 $\Rsim(\G)$ 与目标场尽可能一致，同时满足物理与工艺限制。因此引入如下评分函数：
\[
\score(\G;\Rtar)=
\alpha\,\loss_{\mathrm{resp}}(\Rsim,\Rtar)
+\beta\,\loss_{\mathrm{pol}}
+\gamma\,\loss_{\mathrm{spec}}
+\delta\,\loss_{\mathrm{phys}},
\]
其中主项衡量目标响应与验证响应之间的失配，额外项分别刻画偏振分离、光谱一致性以及物理可实现性。虽然不同实验的权重可以调整，但整体形式保持不变，因为所有目标都在同一响应空间中定义。

同样重要的是，这一表述天然兼容多解逆问题。对于同一个目标响应，往往存在多个几何结构都能达到足够好的效果。因此这里的目标不是找到唯一解，而是描述一个“由目标响应场约束的可行解集”。这也正是下一节采用生成式逆向设计而非单一回归的根本动机。
